% Korean mirror: sections/01-introduction-ko.tex
\section{Introduction}

In a battle royale match, final placement is not a static outcome; it is the
accumulation of survival risk that changes phase by phase inside a shrinking
battlefield. A team's risk of elimination rises and falls with where the
team stands, which opponents are nearby, and how hard the contracting zone
presses on it. Post-match aggregates---total kills, total damage, final
travel distance---flatten this process into post-hoc statistics stripped of
temporal order. The same amount of damage signals different risk in an early
skirmish than in a late fight inside a tight zone, and the same travel
distance means different exposure when the movement is chosen early than
when it is forced late. A match-level label therefore cannot say when a team
became vulnerable or which spatial condition was responsible.

We study this process in \textit{PUBG: Battlegrounds}. PUBG fits the
question structurally. A match places up to 100 players on a large map and
advances through numbered \textit{phases}: in each phase the game announces
the next \textit{safe zone}---the area where teams may remain---and the
surrounding \textit{blue zone} deals growing damage to anyone outside
it~\cite{pubgGameInfo}. The developer, KRAFTON, exposes a per-match
telemetry stream through its official developer API, recording player
positions at regular intervals together with time-ordered combat, zone, and
elimination events~\cite{pubgTelemetryEvents}. Our unit of analysis is the
team, not the player: player-level records are aggregated into team-level
trajectories, exposure, and inter-team relations, and we define a team's
elimination time as the moment its last surviving member is eliminated. The
stream is spatio-temporal by nature---teams are moving objects, the zones
are dynamic regions, and combat and elimination are events anchored to a
time and a place.

Most public benchmarks and prediction studies on PUBG predict final
placement from post-match aggregates~\cite{kagglePubgFinishPlacement,sen2021pubgFinalRank,lee2025pubgRank}.
That problem setting suits match-level outcome prediction; it is simply not
designed for phase-level survival dynamics. Closer to our aim, Sagong et
al.~\cite{sagong2025hunterCamper} analyze PUBG survival phase by phase,
scoring each player's aggressiveness and activity and predicting next-phase
survival from the resulting strategy types. We inherit their phase-level
problem setting but change the unit of analysis and the representation:
teams instead of players, and a phase-aligned spatio-temporal reconstruction
of the full telemetry---team-level features and dynamic team graphs---instead
of a small set of strategy scores. The aim moves beyond classifying
next-phase survival itself, toward analyzing under which spatio-temporal
conditions and inter-team relations a team's elimination risk takes shape.

Concretely, we predict, from what is observable up to the end of each
phase, the probability that a team is eliminated during the next phase. We
call this quantity the \textit{next-phase elimination hazard}; it is the
discrete-time event-history view of
survival~\cite{allison1982discrete,singer1993discrete}, and it replaces the
single match-level label with one risk label per team and phase. Supporting
this task takes more than feature engineering. Features, labels, and
relations must all be derived on shared time and space axes, phase by
phase, without touching future information; we therefore reconstruct the
telemetry into a phase-aligned spatio-temporal database (STDB) whose
objects---moving teams, dynamic zones, spatial events, and phase
intervals---make those derivations systematic and reproducible. Finally, a
team's risk depends not only on its own state but on the configuration of
teams around it: who is near, who attacks whom, and who is pushed into the
same shrinking space. We encode these relations as a dynamic team graph per
phase and test whether they carry signal of their own.

Two research questions organize the study:
\begin{description}
  \item[\textbf{RQ1.}] Can official PUBG telemetry be reconstructed into a phase-aligned STDB from which team trajectories, zone exposure, combat events, and dynamic inter-team relations are derived as team-phase features and dynamic team graphs?
  \item[\textbf{RQ2.}] On this representation, (i) how accurately can the next-phase elimination hazard be predicted from the team-phase features alone, and (ii) do dynamic team graphs add predictive signal beyond models that use only those features?
\end{description}

This study makes three contributions. First, it reformulates battle-royale
survival prediction from final-placement estimation to a team-level,
phase-level elimination hazard. Second, it provides the representation that
makes this task well-posed: a phase-aligned STDB built from 21,824 matches
collected through the official developer API, with leakage-free team-phase
features and dynamic team graphs. Third, it quantifies where the predictive
information lies. Replacing a strategy-score baseline that mirrors the prior
problem setting with the full telemetry-based representation improves the
same learning algorithm from AUC 0.649 to 0.765 in the pooled aligned
evaluation; dynamic team graphs then add a further improvement over models
that use only the team-phase features, an effect confirmed by match-level
paired tests, edge ablations, and
randomization controls. Beyond accuracy, the phase-level view localizes when
and under which spatial conditions a team's risk rises, and can serve as a
starting point for broadcast analysis of at-risk teams refreshed at each
phase boundary, phase-by-phase
feedback in competitive play, and survival-process analysis of live-service
game telemetry.
